\chapter{ТЕХНИЧЕСКАЯ РЕАЛИЗАЦИЯ}

% ─────────────────────────────────────────────────────────────────
\section{Архитектура системы}
% ─────────────────────────────────────────────────────────────────

\subsection{Контекстная диаграмма (C4 Level~1)}

ИИ-методист~--- один из~микросервисов платформы Adapstory,
образовательной LMS с~моделью B2B~SaaS для~корпоративных
заказчиков. На уровне системного контекста (C4~Level~1,
приложение~Г, раздел~Г.1) он взаимодействует со~следующими
акторами и~системами:

\begin{itemize}
  \item \textbf{Разработчик курса / методист / HR-специалист} --- основной пользователь:
    загружает исходные документы, инициирует подготовку черновика,
    просматривает, корректирует и~утверждает результат;
  \item \textbf{BFF~School} (Backend-for-Frontend школы) --- единственная
    точка входа для~веб-клиента; маршрутизирует запросы к~микросервису;
  \item \textbf{Адаптивная LMS Adapstory} --- принимает утверждённую структуру
    курса через REST~API и~публикует её в~каталоге;
  \item \textbf{Хранилище документов MinIO} --- источник исходных файлов
    (.pdf, .docx, .pptx, .xlsx) и~хранилище медиафайлов курса;
  \item \textbf{Ollama GPU-сервер} --- провайдер локальных LLM
    (семейство Qwen~2.5, FLUX.1-dev);
  \item \textbf{Qdrant} --- векторная база знаний для~RAG.
\end{itemize}

\subsection{Диаграмма контейнеров (C4 Level~2)}

На уровне контейнеров (приложение~Г, раздел~Г.2) микросервис
разделён на~процессы:

\begin{itemize}
  \item \textbf{FastAPI~API~Gateway} --- принимает HTTP-запросы от~BFF,
    валидирует входные данные (Pydantic~v2), публикует задания в~Celery;
  \item \textbf{Celery~Workers} (×N) --- пул асинхронных воркеров,
    исполняющих задачи Document~Ingestion и~Course~Generation;
  \item \textbf{LLM~Gateway} --- провайдер-агностичный интерфейс
    к~языковым моделям с~retry, circuit~breaker (Resilience4j-аналог
    на~Python / Tenacity) и~кэшированием промптов;
  \item \textbf{LangGraph~Runtime} --- исполняет граф агентов,
    сохраняет checkpoint в~Redis;
  \item \textbf{Haystack~AsyncPipeline} --- конвейер индексации документов.
\end{itemize}

\subsection{Диаграмма компонентов (C4 Level~3)}

Декомпозиция компонентов LangGraph~Runtime приведена
в~приложении~Г (раздел~Г.3) и~раскрыта в~разделе~3.2.

\subsection{Доменная модель}

Доменная модель ИИ-методиста включает сущности:

\begin{itemize}
  \item \textbf{EducationalProduct} --- запрос на~проектирование курса:
    содержит ссылку на~исходный документ, выбранную методологию
    (ADDIE / SAM / Backward~Design) и~параметры целевой аудитории;
  \item \textbf{CourseBlueprint} --- рабочий методический черновик структуры курса:
    метаданные, список модулей, учебные цели с~разметкой по~Блуму;
  \item \textbf{Module} --- логический раздел курса: название, описание,
    список уроков, competency~tags;
  \item \textbf{Lesson} --- атомарная единица обучения: текст, список
    иллюстраций, ссылка на~TTS-подкаст, тестовые вопросы;
  \item \textbf{CompetencyTag} --- тег компетенции из~графа знаний Neo4j.
\end{itemize}

% ─────────────────────────────────────────────────────────────────
\section{Мультиагентная система (LangGraph)}
% ─────────────────────────────────────────────────────────────────

\subsection{Граф агентов (StateGraph)}

Граф агентов реализован через \texttt{StateGraph}
из~LangGraph~\cite{langgraph_docs_2025}.
Состояние графа (\texttt{CourseState}) хранит входной документ,
промежуточные артефакты (структуру курса, список модулей, контент уроков)
и~флаги человеческого ревью. Полная спецификация графа~--- в~приложении~Д.

\subsection{Агенты конвейера}

Конвейер включает пять агентов:

\begin{enumerate}
  \item \textbf{TargetAnalysisAgent} --- анализирует исходный документ,
    определяет целевую аудиторию, предварительные знания и~разрыв
    между текущим и~желаемым уровнем компетентности;
  \item \textbf{CourseStructureAgent} --- формирует структуру курса
    (CourseBlueprint) в~соответствии с~выбранной методологией (ADDIE / SAM / BD)
    и~формирует учебные цели с~разметкой по~таксономии Блума;
  \item \textbf{ContentCreatorAgent} --- подготавливает черновые тексты
    уроков с~использованием RAG: извлекает релевантные фрагменты
    из~Qdrant и~передаёт их в~контекст LLM;
  \item \textbf{FactCheckAgent} --- верифицирует фактические утверждения
    в~сгенерированном контенте относительно исходного документа,
    выявляет галлюцинации и~инициирует пересмотр при~низкой достоверности
    относительно источника (faithfulness);
  \item \textbf{MediaGenerationAgent} --- параллельно подготавливает
    иллюстрации (FLUX.1-dev) и~TTS-подкасты (Silero) для~уроков.
\end{enumerate}

\subsection{Контур участия человека}

После~завершения CourseStructureAgent граф приостанавливается
на~ноде \texttt{human\_review\_structure} через механизм
\texttt{interrupt\_before}~\cite{langgraph_docs_2025}.
Разработчик курса получает уведомление (веб-сокет / email) и~может:
одобрить структуру, отредактировать отдельные модули или~запросить
полную перегенерацию с~альтернативной методологией.
После~одобрения граф возобновляется с~сохранённого checkpoint.

\subsection{Управление состоянием и трассировка}

Checkpoint~API LangGraph сохраняет \texttt{CourseState}
после каждой ноды в~Redis, что даёт:

\begin{itemize}
  \item идемпотентное возобновление при~сбое воркера;
  \item возможность просмотра истории выполнения через LangSmith;
  \item параллельное исполнение ветки MediaGenerationAgent без~блокировки
    основного конвейера.
\end{itemize}

LangSmith~\cite{langsmith_docs_2025} трассирует
каждый вызов LLM: входной промпт, ответ модели, latency, токены~---
достаточно для~анализа деградации качества на~конкретных документах
без~воспроизведения всего конвейера.

% ─────────────────────────────────────────────────────────────────
\section{Ключевые конвейеры}
% ─────────────────────────────────────────────────────────────────

\subsection{Конвейер обработки документов (Document Ingestion Pipeline, Haystack~2.x)}

Индексация документов построена на~Haystack~2.x
\texttt{AsyncPipeline}~\cite{haystack2_docs_2024}:

\begin{enumerate}
  \item \textbf{MinIODocumentFetcher} --- загружает файл из~MinIO
    по~presigned~URL (форматы: PDF, DOCX, PPTX, XLSX);
  \item \textbf{DocumentConverter} --- конвертирует файл в~текст
    (pdfminer / python-docx / python-pptx); сохраняет метаданные источника;
  \item \textbf{SemanticDoubleMergingSplitter} --- разбивает текст
    на~семантические фрагменты;
  \item \textbf{OllamaTextEmbedder} --- векторизует фрагменты
    через модель \texttt{deepvk/USER-bge-m3} (Ollama API);
  \item \textbf{QdrantDocumentStore.write\_documents()} --- записывает
    фрагменты с~векторами в~Qdrant (коллекция курса).
\end{enumerate}

Celery-воркер запускает конвейер асинхронно сразу после загрузки
документа; генерация курса начинается только по~завершении индексации.

\subsection{Конвейер подготовки черновика курса}

Черновик курса формируется последовательностью вызовов агентов
LangGraph с~общим \texttt{CourseState}:

\begin{enumerate}
  \item TargetAnalysisAgent формирует профиль аудитории;
  \item CourseStructureAgent выбирает промпт-шаблон
    (ADDIE / SAM / Backward~Design, приложение~Е) и~формирует
    структурированный JSON (Pydantic~\texttt{CourseBlueprint});
  \item ContentCreatorAgent для~каждого урока выполняет гибридный поиск
    в~Qdrant (dense~+~sparse), формирует контекст и~подготавливает черновой текст урока;
  \item FactCheckAgent проверяет каждое утверждение в~тексте урока
    против~исходного документа;
  \item ReviewAgent формирует финальный отчёт о~качестве курса.
\end{enumerate}

\subsection{Мультимедиа-конвейер}

После~одобрения текстовой структуры MediaGenerationAgent
запускается параллельно с~ContentCreatorAgent через \texttt{Send~API}
LangGraph. Для~каждого урока:

\begin{itemize}
  \item подготавливается иллюстрация: заголовок урока + ключевые термины
    $\rightarrow$ FLUX.1-dev~\cite{flux_2024} $\rightarrow$ PNG $\rightarrow$ MinIO;
  \item подготавливается аудиоподкаст: LLM~формирует диалог двух ведущих
    (Наставник + Ученик) по~тексту урока $\rightarrow$ Silero~TTS~v4~\cite{silero_tts_2024}
    $\rightarrow$ MP3 $\rightarrow$ MinIO.
\end{itemize}

\subsection{Механизм актуализации}

Изменение исходного документа запускает конвейер актуализации:

\begin{enumerate}
  \item Document~Ingestion повторно индексирует обновлённый документ;
  \item FactCheckAgent прогоняет все существующие уроки против
    нового индекса;
  \item Уроки с~faithfulness ниже порога (по~умолчанию 0.8) помечаются
    как~\texttt{OUTDATED} и~передаются ContentCreatorAgent
    на~перегенерацию;
  \item контур участия человека уведомляет владельца курса о~изменённых разделах.
\end{enumerate}

% ─────────────────────────────────────────────────────────────────
\section{Инфраструктура и развёртывание}
% ─────────────────────────────────────────────────────────────────

\subsection{Dev-стенд: K3s + GPU}

Сервис развёрнут в~K3s-кластере Adapstory
(single-node, 192.168.1.3) под~управлением ArgoCD~(App-of-Apps pattern).
Конфигурация оформлена как Helm-чарт в~GitOps-репозитории~\cite{adapstory_gitops_2025}.

GPU-узел (NVIDIA~Tesla~T4) предоставлен через NVIDIA~GPU~Operator
с~time-slicing: Ollama и~Celery-воркеры делят GPU-ресурс.
Все~персональные данные и~документы клиентов остаются внутри периметра
локального развёртывания, что соответствует требованиям
152-ФЗ~\cite{fz_152_2006}.

\subsection{LLM~Gateway}

LLM~Gateway предоставляет единый асинхронный
провайдер-агностичный API (\texttt{complete(prompt, model, params)}):

\begin{itemize}
  \item \textbf{Retry} (Tenacity): экспоненциальный backoff,
    до~3~попыток при~HTTP~5xx от~Ollama;
  \item \textbf{Circuit~Breaker}: при~недоступности Ollama >~30~с
    переключается на~резервный провайдер (OpenAI~API);
  \item \textbf{Кэш промптов}: идентичные промпты кэшируются в~Redis
    (TTL~1~час) для~снижения нагрузки на~GPU.
\end{itemize}

\subsection{CI/CD и мониторинг}

Jenkins-пайплайн собирает образы: Maven/uv~build~$\rightarrow$
Docker build (Kaniko, linux/amd64)~$\rightarrow$ push в~Harbor-реестр~$\rightarrow$
обновление тега в~GitOps-репозитории~$\rightarrow$ ArgoCD~sync.

OpenTelemetry~Collector собирает метрики:
latency по~нодам LangGraph-графа, throughput Celery-очереди,
GPU-utilization. Дашборды доступны в~Grafana (dev-стенд).
